\scanpage{101}
% archive sheet number: 100
\textcircled{1} $X$ schéma
\[
\begin{aligned}
D_i&\quad\text{diviseurs de Cartier }\geq0 &&1\leq i\leq\unsure{n|r}\\
n_i&\quad\text{entiers }\geq1 &&1\leq i\leq r
\end{aligned}
\]
$\Gamma$ schéma en groupes sur $S$

$\deleted{\unsure{\mu=}}\mu_{\underline n}=\prod\mu_{n_i}\qquad(\underline n=(n_1,-,n_r))$ sur $X$

Extension de groupes sur $X$ \quad $1\to\mu\mathbin{-}G\mathbin{-}\Gamma\mathbin{-}1$

$X'\longrightarrow\widetilde X_1\longrightarrow X$ morphismes de schémas sur $X$, à op. $G$.

On suppose a) que $G$ opère sur $\widetilde X_1$ que $\mu$ y opère
trivialement, que \deleted{\unsure{le}} $\widetilde X_1$ devient \uline{torseur} sous $\Gamma$

b) En tant que fibré à opérateurs $\mu_{\widetilde X}$\inserted{$=\mu_{\underline n,\widetilde X}$}, $X'/$ sur $X$,
est loc. \deleted{\unsure{à|$\mathcal U$}} (sur $\widetilde X_1$) de type kummerien
relatif aux images inverses $\widetilde D_i$ des $D_i$
[\unsure{supposées|supposés} rester \unsure{encore|comme} diviseurs) \inserted{d'exposant $n_i$}, i.e. de la forme
% margin note written sideways
\margin{\diagram{\begin{tikzcd}[ampersand replacement=\&,column sep=tiny,row sep=small,cramped]
B_G\arrow[d] \& \unsure{B|\theta}_{\mu}\arrow[l]\arrow[d] \& B_G/(\widetilde X,G)\arrow[dl]\arrow[d,no head]\\
B_\Gamma \& (\widetilde X,\Gamma)\arrow[l,Rightarrow]\arrow[r,phantom,"\simeq"] \& X
\end{tikzcd}}}
\[
\widetilde X[T_1,..,T_r]/(T_1^{n_1}-\deleted{\unsure{x}}a_1,..,T_r^{n_r}-a_r)
\]
, où les $a_i$ sont des équations des $\widetilde D_i$ (locales).

\deleted{b) $\Gamma$ opérant sur $\mu$, définit par} On se donne un isom $\widetilde X\mathbin{\overset{\Gamma}{\times}}\mu\simeq\mu_{\underline n}$

L'extension \inserted{$G$} de $\Gamma$ par $\mu$ définit un élément
\[
\deleted{\unsure{a}}\alpha\in H^2(B_\Gamma,\mu)\deleted{\unsure{\simeq\prod H^2(B_\Gamma,\mu_{n_i})}} .
\]
Or on a un hom canonique, qui au \inserted{$\Gamma$-torseur} $X$
\diagram{\begin{tikzcd}[ampersand replacement=\&,column sep=small,row sep=tiny]
H^2(B_\Gamma,\mu)\arrow[r,"\varphi"] \& H^2(X,\mu_{\underline n})\arrow[d,phantom,"{\rotatebox{90}{\(\simeq\)}}"] \& {}\simeq\prod H^{\unsure{1|2}}(X,\mu_{n_i})\\
\& H^2(B_\Gamma/(\widetilde X,\Gamma),\mu) \&
\end{tikzcd}}

On veut prouver que, si
\[
\varphi(\alpha)=(\xi_i)_{1\leq i\leq r}\qquad \xi_i=c\inserted{\substack{\text{1\textsuperscript{re} classe}\\\text{de Chern}\\\text{mod }n_i}}(\underline{\mathrm O}_X(D_i)\deleted{\unsure{\text{fibré inv}}})\pmod{n_i}.
\]

\textcircled{2} Variante schémas formels, (ou topos loc. annelés généraux)
